\part{Aktueller Stand}

\chapter{Woran es liegt}

\section{Die Tabelle aus Kapitel 1, ein letztes Mal}

\begin{merksatz}
\boxtitel{Was inzwischen dazugehört}
Am Anfang stand eine Tabelle, die zeigte, dass früh behandelte Patienten häufiger starben.

\begin{center}
{\small
\renewcommand{\arraystretch}{1.5}
\begin{tabular}{@{}P{1.2cm}P{3.6cm}P{3.6cm}@{}}
\toprule
\textbf{Kap.} & \textbf{Frage} & \textbf{Antwort}\\
\midrule
$1$ & Was zeigt die Tabelle? & Nichts Kausales.\\
$2$ & Was will man wissen? & $\mathbb{E}[Y(1)-Y(0)]$.\\
$3$ & Wie käme man hin? & Auslosen.\\
$4$--$6$ & Und ohne Los? & Hinterwege sperren.\\
$7$--$8$ & Womit? & Neigung, Gewichte.\\
$9$--$11$ & Wenn das nicht geht? & Instrument, Sprung, Nachbildung.\\
$12$--$14$ & Über die Zeit, für wen? & g-Formel, $\tau(\vek z)$.\\
\bottomrule
\end{tabular}}
\end{center}

\begin{center}
\begin{minipage}{0.9\textwidth}
\textbf{Keine dieser Antworten steckt in den Daten. Sie alle bestehen darin, die richtige Annahme zu benennen und ihre Folgen auszurechnen.}
\end{minipage}
\end{center}
\end{merksatz}

\section{Vier Muster}

\begin{merksatz}
\boxtitel{Muster 1 --- die Daten entscheiden nicht}
\begin{center}
{\small
\renewcommand{\arraystretch}{1.48}
\begin{tabular}{@{}P{1.4cm}P{7.0cm}@{}}
\toprule
$1$ & Satz 1.2: Dieselbe Tabelle, Effekt $+0{,}60$ oder null.\\
$4$ & Drei Graphen sind ununterscheidbar; nur der Kollider verrät sich.\\
$6$ & Manchmal gibt es gar keine zulässige Anpassungsmenge.\\
$13$ & Die Kreuzwelt-Annahme ist grundsätzlich unprüfbar.\\
\bottomrule
\end{tabular}}
\end{center}

\begin{center}
\begin{minipage}{0.88\textwidth}
\kursiv{\textbf{Das ist keine Schwäche des Faches, sondern sein wichtigster Beitrag.} Vor diesen Verfahren rechnete man Regressionen und sprach von \enquote{Einfluss}; die Annahmen blieben unausgesprochen und damit unbestreitbar. Wer heute einen Graphen zeichnet, legt sie offen --- und jeder Fachkollege kann einen Pfeil bestreiten.}
\end{minipage}
\end{center}
\end{merksatz}

\begin{merksatz}
\boxtitel{Muster 2 --- mehr anpassen ist nicht besser}
\begin{center}
{\small
\renewcommand{\arraystretch}{1.48}
\begin{tabular}{@{}P{1.4cm}P{7.0cm}@{}}
\toprule
$5$ & Mediator: Effekt von $0{,}72$ auf $0{,}004$.\\
$5$ & Kollider: aus null wird $-0{,}39$.\\
$5$ & Instrument: Verzerrung verdoppelt.\\
$12$ & Dieselbe Größe muss zugleich angepasst und durchlaufen werden.\\
\bottomrule
\end{tabular}}
\end{center}

\begin{center}
\begin{minipage}{0.88\textwidth}
\kursiv{\textbf{In jedem dieser Fälle sieht das falsche Modell besser aus als das richtige} --- es erklärt mehr Streuung, hat engere Intervalle und bestandene Gütemaße. Kein Auswahlverfahren aus Band 22 findet den Fehler, weil er nicht in den Daten steckt.}
\end{minipage}
\end{center}
\end{merksatz}

\begin{merksatz}
\boxtitel{Muster 3 --- der Entwurf schlägt das Verfahren}
\begin{center}
{\small
\renewcommand{\arraystretch}{1.48}
\begin{tabular}{@{}P{1.4cm}P{7.0cm}@{}}
\toprule
$3$ & Nach Behandlung statt nach Zuteilung: fast doppelter Effekt.\\
$11$ & Unsterbliche Zeit: $73$ Prozent Gewinn aus nichts.\\
$11$ & Die Hormonstudien widersprachen sich wegen des Beginns, nicht wegen der Störgrößen.\\
$14$ & Untergruppen brauchen die achtfache Fallzahl.\\
\bottomrule
\end{tabular}}
\end{center}

\begin{center}
\begin{minipage}{0.88\textwidth}
\kursiv{\textbf{Die größten Fehler dieses Bandes sind Entwurfsfehler, keine Rechenfehler} --- und sie sind vermeidbar. Das ist dieselbe Lehre wie in Band 22, Kapitel~3, wo die Prüfung wichtiger war als das Verfahren. Wer die sieben Zeilen der Zielstudie ausfüllt, bevor er rechnet, hat mehr gewonnen als durch jede Verfeinerung des Schätzers.}
\end{minipage}
\end{center}
\end{merksatz}

\begin{merksatz}
\boxtitel{Muster 4 --- jede Zahl gilt für eine Bevölkerung}
\begin{center}
{\small
\renewcommand{\arraystretch}{1.48}
\begin{tabular}{@{}P{1.4cm}P{7.0cm}@{}}
\toprule
$2$ & $\tau$ und $\tau_{\text{beh}}$: $-0{,}375$ gegen $-0{,}250$.\\
$9$ & Der Instrumentschätzer gilt für die Befolger --- die niemand benennen kann.\\
$10$ & Der Sprung gilt an der Schwelle.\\
$14$ & Der Durchschnitt von $-0{,}10$ gilt für niemanden.\\
\bottomrule
\end{tabular}}
\end{center}

\begin{center}
\begin{minipage}{0.88\textwidth}
\kursiv{\textbf{Ein kausaler Effekt ohne Angabe der Bevölkerung ist unvollständig} --- so unvollständig wie ein Mittelwert ohne Angabe, worüber gemittelt wurde. Die Verfahren dieses Bandes unterscheiden sich weniger in ihrer Güte als in der Frage, die sie beantworten.}
\end{minipage}
\end{center}
\end{merksatz}

\section{Was offen ist}

\begin{merksatz}
\boxtitel{Vier Fragen}
\begin{center}
{\small
\renewcommand{\arraystretch}{1.55}
\begin{tabular}{@{}P{3.0cm}P{5.6cm}@{}}
\toprule
\textbf{Übertragung} & Wann gilt ein Effekt in einer anderen Klinik? Kapitel~6 gibt ein Kriterium --- es verlangt zu wissen, worin sich die Häuser unterscheiden.\\
\textbf{Struktur finden} & Graphen aus Daten zu gewinnen, gelingt nur bis zur Äquivalenzklasse.\\
\textbf{Entscheiden} & Strategien lernen statt bewerten --- und die Positivitätsgrenze.\\
\textbf{Verbinden} & Vorhersage und Ursache in einem Modell.\\
\bottomrule
\end{tabular}}
\end{center}

\begin{center}
\begin{minipage}{0.88\textwidth}
\kursiv{Die letzte Frage schließt den Bogen zu Band 22. \textbf{Ein Modell, das nur ursächliche Größen benutzt, bleibt bei Verteilungsverschiebungen gültig und sagt schlechter vorher; eines, das alles benutzt, sagt besser vorher und bricht bei Verschiebungen zusammen.} Das Sepsis-Warnsystem von 2021 war vom zweiten Typ. Ob es einen brauchbaren Mittelweg gibt, ist offen.}
\end{minipage}
\end{center}
\end{merksatz}

\section{Sechs Sätze, die bleiben}

\begin{merksatz}
\boxtitel{Die Bilanz}
\begin{center}
{\small
\renewcommand{\arraystretch}{1.55}
\begin{tabular}{@{}P{0.7cm}P{7.7cm}@{}}
\toprule
$1$ & Zusammenhang entsteht aus Wirkung, gemeinsamer Ursache oder Auswahl --- und die Daten sagen nicht, woraus.\\
$2$ & Der individuelle Effekt ist unbeobachtbar; der durchschnittliche ist unter angebbaren Annahmen schätzbar.\\
$3$ & Das Los stellt her, was sonst behauptet werden muss.\\
$4$ & Eine Menge ist zulässig, wenn sie keinen Nachkommen enthält und jeden Hinterweg sperrt.\\
$5$ & Alle Störgrößen lassen sich durch eine einzige Zahl ersetzen --- die Behandlungsneigung.\\
$6$ & Wird über die Zeit behandelt, ist dieselbe Größe Störgröße und Mediator zugleich; nur die g-Formel trennt beides.\\
\bottomrule
\end{tabular}}
\end{center}
\end{merksatz}

\begin{merksatz}
\boxtitel{Ein Schlusswort}
\begin{center}
\begin{minipage}{0.9\textwidth}
Dieser Band begann mit einer Frage, die jeder Arzt täglich stellt: Was geschieht, wenn ich dies tue statt jenes?

Fünfzehn Kapitel später lässt sich sagen, was die Mathematik darauf antwortet. Sie antwortet nicht mit einer Zahl, sondern mit einer Buchführung: Welche Annahme brauche ich, damit diese Rechnung gilt? Welchen Weg muss ich sperren, welchen offen lassen? Für welche Patienten gilt das Ergebnis? Und woran würde ich merken, dass ich falsch liege?

Die Antworten sind unbequem. Fast jede wichtige Annahme ist unprüfbar. Fast jede Zahl gilt für eine andere Bevölkerung als die, für die man sie braucht. Und die größten Fehler entstehen nicht beim Rechnen, sondern beim Entwerfen --- lange bevor der erste Schätzer gewählt ist.

Aber die Buchführung ist mehr, als es vorher gab. Wer die Annahmen aufschreibt, kann sie bestreiten lassen. Wer die Bevölkerung nennt, kann gefragt werden, ob sie die richtige ist. Das ist keine Gewissheit --- es ist die Form, in der man über Ursachen streiten kann, ohne aneinander vorbeizureden.
\end{minipage}
\end{center}

\begin{center}
\begin{minipage}{0.88\textwidth}
\kursiv{Damit steht dieser Band am Ende der zweiten Werkhälfte. Die Statistik sagt, was die Daten hergeben. Die Optimierung sagt, was es kostet. Die Informationstheorie sagt, was übertragbar ist. Die Spieltheorie sagt, was übrig bleibt, wenn der andere mitdenkt. Das Lernen sagt, was auf den nächsten Fall zutrifft. \textbf{Und die Kausalität sagt, was davon sich ändert, wenn man handelt.}}
\end{minipage}
\end{center}
\end{merksatz}
